%%%%%%%%%%%%%%%%%%%%%%%%%%%%%%%%%%%%%%%%%%%%%%%%%%%%%%%%%%%%%%%%%%%%%%
%% Copyright (c) 2016 Carsten Wulff Software, Norway
%% %%%%%%%%%%%%%%%%%%%%%%%%%%%%%%%%%%%%%%%%%%%%%%%%%%%%%%%%%%%%%%%%%%%
%% Created       : wulff at 2016-11-5
%% %%%%%%%%%%%%%%%%%%%%%%%%%%%%%%%%%%%%%%%%%%%%%%%%%%%%%%%%%%%%%%%%%%%
%% This program is free software: you can redistribute it and/or modify
%% it under the terms of the GNU General Public License as published by
%% the Free Software Foundation, either version 3 of the License, or
%% (at your option) any later version.
%%
%% This program is distributed in the hope that it will be useful,
%% but WITHOUT ANY WARRANTY; without even the implied warranty of
%% MERCHANTABILITY or FITNESS FOR A PARTICULAR PURPOSE.  See the
%% GNU General Public License for more details.
%%
%% You should have received a copy of the GNU General Public License
%% along with this program.  If not, see <http://www.gnu.org/licenses/>.
%%%%%%%%%%%%%%%%%%%%%%%%%%%%%%%%%%%%%%%%%%%%%%%%%%%%%%%%%%%%%%%%%%%%%%
%\documentclass[final,11pt]{IEEEtran}
\documentclass[crop,border=0pt]{standalone}
\usepackage{times}


%\usepackage{placeins}
\usepackage{upgreek}
\usepackage{amssymb,amsmath}
\usepackage{cite}
\usepackage{verbatim}
\usepackage{listings}
\usepackage{xcolor}
\usepackage{flafter}
\usepackage{booktabs}
\usepackage{url}
\usepackage{ textcomp }

\hyphenation{op-tical net-works semi-conduc-tor}


% -JSON styel

%\definecolor{darkGreen}{rgb}{0,0.2097,0}
%\newcommand{\myfigwidth}{\linewidth}
%\newcommand{\myfigwidtha}{\linewidth}
%\newcommand{\myfigwidthl}{\textwidth}
%\newcommand{\myfigwidthc}{0.6\linewidth}


\newcommand{\myfigwidtha}{3.3in}
\newcommand{\myfigwidthl}{7in}
\newcommand{\myfigwidthc}{2in}
%\newcommand{\myfigwidthe}{3.2in}
%\newcommand{\myfigwidthb}{5in}

\newcommand{\myfigname}{fig_}
\newcommand{\reg}[2]{{\figurename} \ref{#1}(#2)}
\newcommand{\req}[1]{{\figurename} \ref{#1}}
%\newcommand{\myeqname}{eq:}
%\newcommand{\req}[1]{(\ref{\myeqname#1})}

\definecolor{armygreen}{rgb}{0, 0.5, 0}
\newcommand{\missing}[1]{{ #1}}
\newcommand{\edit}[1]{{ #1}}
\newcommand{\secedit}[1]{{ #1}}
\newcommand{\deleted}[1]{{ #1}}
\newcommand{\moved}[1]{{ #1}}

%\newcommand{\missing}[1]{{\color{red} #1}}
%\newcommand{\edit}[1]{{\color{red} #1}}
%\newcommand{\secedit}[1]{{\color{armygreen} #1}}
%\newcommand{\deleted}[1]{{\color{violet} #1}}
%\newcommand{\moved}[1]{{\color{blue} #1}}

\newcommand{\SD}{Sigma-Delta }
\newcommand{\eqn}[1]{
  \begin{math}
    #1
  \end{math}}

\newcommand\JSONnumbervaluestyle{\color{blue}}
\newcommand\JSONstringvaluestyle{\color{red}}



\newif\ifcolonfoundonthisline

\makeatletter

\lstdefinestyle{json}
{
  showstringspaces    = false,
  keywords            = {false,true},
  alsoletter          = 0123456789.,
  morestring          = [s]{"}{"},
  stringstyle         = \ifcolonfoundonthisline\JSONstringvaluestyle\fi,
  MoreSelectCharTable =%
  \lst@DefSaveDef{`:}\colon@json{\processColon@json},
  basicstyle          = \fontsize{7}{7}\ttfamily,
  keywordstyle        = \fontsize{7}{7}\ttfamily\bfseries,
}

% flip the switch if a colon is found in Pmode
\newcommand\processColon@json{%
  \colon@json%
  \ifnum\lst@mode=\lst@Pmode%
  \global\colonfoundonthislinetrue%
  \fi
}

\lst@AddToHook{Output}{%
  \ifcolonfoundonthisline%
  \ifnum\lst@mode=\lst@Pmode%
  \def\lst@thestyle{\JSONnumbervaluestyle}%
  \fi
  \fi
  % override by keyword style if a keyword is detected!
  \lsthk@DetectKeywords%
}

% reset the switch at the end of line
\lst@AddToHook{EOL}%
{\global\colonfoundonthislinefalse}

\makeatother

% \lstset{  basicstyle=\tiny}

% \lstset{
% basicstyle=\fontsize{11}{13}\selectfont\ttfamily
% }


\lstset{language=Java}




%\usepackage[paperwidth=\maxdimen,paperheight=\maxdimen]{geometry}


\usepackage{tikz,pgfplots}
\usepackage{ifthen}
\usepackage{calc}
\usetikzlibrary{patterns}
%\usepackage[tightpage,active]{preview}
\usetikzlibrary{circuits.logic.US}
\usetikzlibrary{circuits.ee.IEC}
\usepgflibrary{shapes.arrows}
\usepgflibrary{arrows}
\usepackage{circuitikz}


% circuitikz draws bipole outlines (resistors, diodes, capacitors) at
% bipoles/thickness times the ambient line width, and the default is 2.
% That makes a circuitikz resistor visibly heavier than the wire it sits
% on, and heavier than the hand rolled zig-zag in ckt_lib.tex. One is the
% house width: everything is drawn at the environment's own thickness.
% Arrow tips: one filled tip everywhere, set through the '>' knob so a
% figure only ever writes ->, <- or <->. Latex (from arrows.meta)
% scales with the line width, so it stays in proportion if a drawing
% is ever set thicker or thinner than the house 'thick'.
\usetikzlibrary{arrows.meta}
\tikzset{>={Latex}}

\ctikzset{bipoles/thickness=1}
\ctikzset{tripoles/mos style/arrows}
\ctikzset{tripoles/pmos style/emptycircle}
%\PreviewEnvironment{circuitikz}
\pagestyle{empty}


\tikzstyle{cicfill}=[minimum size=1.4em]
\tikzstyle{stuff_fill}=[rectangle,fill=white,minimum size=1.4em]

\newcommand{\grayOne}{0.3 }
\newcommand{\grayTwo}{0.15 }
\newcommand{\grayThree}{0.26 }
\newcommand{\grayFour}{0.41 }
\newcommand{\grayFive}{0.61 }
\newcommand{\graySix}{0.85 }

%- Gray tone
%\definecolor{active}{rgb}{\graySix,\graySix,\graySix}
%\definecolor{poly}{rgb}{\grayFive,\grayFive,\grayFive}
%\definecolor{contact}{rgb}{\graySix,\graySix,\graySix}
%\definecolor{mOne}{rgb}{\grayFour,\grayFour,\grayFour}
%\definecolor{mTwo}{rgb}{\graySix,\graySix,\graySix}
%\definecolor{mThree}{rgb}{\grayThree,\grayThree,\grayThree}
%\definecolor{mFour}{rgb}{\grayTwo,\grayTwo,\grayTwo}

\definecolor{grey}{rgb}{1,1,1}
\definecolor{lightgrey}{rgb}{\grayFour,\grayFour,\grayFour}
\definecolor{active}{rgb}{0.4,1,0.6}
\definecolor{poly}{rgb}{1.0,0.5,0.5}
\definecolor{cut}{rgb}{0.91,0.91,0}
\definecolor{mOne}{rgb}{0.36,0.36,1}
\definecolor{mTwo}{rgb}{0.655,0.655,0}
\definecolor{mThree}{rgb}{0.3,1,1}
\definecolor{mFour}{rgb}{0,0.2097,0}
\begin{document}


% Why a trap a few nanometres into the oxide has a time constant
% anywhere from microseconds to months.
%
% Upper panel: the band diagram through gate, oxide and silicon in
% strong inversion. The inversion electrons sit in the lowest subband
% E_0 of the surface well (as in mos_2deg). Traps in the oxide with a
% level near E_F can swap an electron with the channel.
%
% Lower panel: the channel wave function's probability density on a
% log axis. In the silicon it is the familiar hump; in the oxide it is
% an evanescent tail, exp(-2 kappa x), a straight line on a log axis.
% With a 3.1 eV barrier and m_ox = 0.5 m_0, kappa = 6.4 /nm, so the
% density falls a decade every 0.18 nm. The tunnelling rate goes as
% the density at the trap, so tau = tau_0 exp(2 kappa x), with
% tau_0 ~ 0.1 ns: 40 us at 1 nm, 13 s at 2 nm, two months at 3 nm.
%
% The energy scale of the upper panel is a cartoon; the lower panel's
% slope in the oxide is computed (0.2 units per decade, 1.2 units per
% nm), and the silicon hump is drawn, not computed.

\begin{tikzpicture}[thick]

  \def\nm{1.2}       % one nanometre, horizontally
  \def\xg{-4.6}      % left edge of the gate
  \def\xox{-3.6}     % gate/oxide interface (3 nm oxide)
  \def\xmax{6.6}     % into the silicon

  % ================= upper panel: the bands =================
  \begin{scope}[shift={(0,1.4)}]
    % gate metal, filled to its Fermi level
    \fill[gray!25] (\xg,-0.3) rectangle (\xox,0.4);
    \draw[gray!60] (\xg,-0.3) rectangle (\xox,4.8);
    \draw[dashed] (\xg,0.4) -- (\xox,0.4);
    \node[anchor=south, scale=0.85] at (-4.1,0.42) {$E_{F,gate}$};
    \node[gray!50!black, anchor=center, scale=0.9] at (-4.1,3.2) {gate};

    % the oxide: a tall barrier, tilted by the gate field
    \fill[oxteal!45] (\xox,-0.3) rectangle (0,4.8);
    \draw[gray!60] (\xox,-0.3) rectangle (0,4.8);
    \draw[blue, very thick] (\xox,3.7) -- (0,4.4);
    \node[oxteal!45!black, anchor=south, scale=0.9] at (-1.8,4.42) {oxide};
    \draw[gray!70, <->] (0.15,0.6) -- (0.15,4.4);
    \node[gray!50!black, anchor=west, scale=0.8] at (0.3,3.3) {3.1 eV};

    % the silicon conduction band: a well at the surface
    \draw[blue, very thick] (0,4.4) -- (0,0.6)
      .. controls (1.6,1.55) and (3.4,2.2) .. (\xmax,2.4);
    \node[blue, anchor=south east, scale=0.95] at (\xmax,2.42) {$E_C$};

    % the lowest subband and the Fermi level
    \draw[armygreen, dashed, very thick] (0,0.95) -- (2.9,0.95);
    \node[armygreen, anchor=west, scale=0.9] at (2.95,0.95) {$E_0$};
    \draw[dashed] (\xox,1.25) -- (\xmax,1.25);
    \node[anchor=south east, scale=0.85] at (\xmax,1.27) {$E_F$};

    % traps near E_F, at one and two nanometres
    \draw[red, line width=2.2pt] ({-1*\nm-0.18},1.38) -- ({-1*\nm+0.18},1.38);
    \draw[red, line width=2.2pt] ({-2*\nm-0.18},1.12) -- ({-2*\nm+0.18},1.12);
    \node[red, anchor=south, scale=0.85] at ({-1.5*\nm},1.75) {traps near $E_F$};

    % the electron tunnels from the channel into the trap
    \draw[red, very thick, ->] (0.55,0.95) .. controls (0.1,1.05) and (-0.4,1.3) .. ({-1*\nm+0.22},1.38);
    \node[red, anchor=north, scale=0.85] at (-0.45,0.85) {tunnel};
  \end{scope}

  % the two trap depths, carried down into the lower panel
  \foreach \d in {1,2} {
    \draw[red!50, thin, dashed] ({-\d*\nm},2.5) -- ({-\d*\nm},{-1.11*\d});
  }

  % ================= lower panel: log of the density =================
  % 0.2 units per decade; the interface value is the reference (0)
  \draw[gray!60] (\xox,-3.7) rectangle (0,0.8);
  \fill[oxteal!20] (\xox,-3.7) rectangle (0,0.8);
  \draw[line width=1.4pt] (0,-3.7) -- (0,0.8);

  % the evanescent tail: -5.56 decades per nm, x in drawing units
  \draw[blue, very thick] plot[domain=\xox:0, samples=2] (\x, {0.2*5.56*\x/\nm});
  % the hump in the silicon, drawn
  \draw[blue, very thick]
    plot[domain=0:\xmax, samples=80]
      (\x, {0.2*ln(1 + 330*(\x/\nm)^2*exp(-1.33*\x/\nm))/ln(10)});
  \node[blue, anchor=south, scale=0.95] at (2.0,0.45) {$\log|\psi_0|^2$};

  \node[blue!60!black, anchor=north west, scale=0.85, align=left] at (-3.5,0.7)
    {straight on a log axis:\\a decade every 0.18 nm};

  % tau at each depth
  \foreach \d in {1,2,3} { \fill[red] ({-\d*\nm},{-1.112*\d}) circle (0.075); }
  \node[red, anchor=west, scale=0.85] at ({-1*\nm+0.1},-1.2) {$\tau \approx 40\,\mu$s};
  \node[red, anchor=west, scale=0.85] at ({-2*\nm+0.1},-2.3) {$\tau \approx 13$ s};
  \node[red, anchor=west, scale=0.85] at ({-3*\nm+0.1},-3.4) {$\tau \approx$ 2 months};

  % depth axis, both ways from the interface
  \draw[gray!60] (\xox,-3.7) -- (\xmax,-3.7);
  \foreach \k in {1,2,3} {
    \draw[gray!60] ({-\k*\nm},-3.7) -- ({-\k*\nm},-3.84);
    \node[gray!50!black, anchor=north, scale=0.8] at ({-\k*\nm},-3.86) {\k};
  }
  \foreach \k in {1,2,3,4,5} {
    \draw[gray!60] ({\k*\nm},-3.7) -- ({\k*\nm},-3.84);
    \node[gray!50!black, anchor=north, scale=0.8] at ({\k*\nm},-3.86) {\k};
  }
  \node[gray!50!black, anchor=north, scale=0.85] at (-1.8,-4.25) {into the oxide [nm]};
  \node[gray!50!black, anchor=north, scale=0.85] at (3.3,-4.25) {into the silicon [nm]};

\end{tikzpicture}
\end{document}
